\section{Introduction}

A basic invariant of a compact Riemannian manifold is its volume. Remarkably, certain conformally compact
asymptotically hyperbolic metrics, including Poincar\'e--Einstein metrics in even dimensions, have a well-defined,
finite renormalized volume. In this paper we study the behavior of this quantity under normalized Ricci flow.

Conformally compact asymptotically hyperbolic (AH) spaces $(M^n, g)$ are a class of complete Riemannian manifolds modeled on the Poincar\'e disk model of hyperbolic space. Any such metric is defined and complete on the interior of a~compact manifold with boundary $\Mbar = M \cup \partial M$, and takes the form $g = \rho^{-2} \gbar$ where $\gbar$ is some smooth metric on $\Mbar$ and where $\rho$ is a smooth boundary defining function. We require that $|{\rm d}\rho/\rho|^2_g \to 1$ at $\partial M$, which ensures that the sectional curvatures of $g$ tend to $-1$ near $\partial M$. The conformal class $[\rho^2 g |_{T\partial M}]$ is a well-defined conformal class on~$\partial M$, and $\del M$ with this conformal class is called the conformal infinity of~$(M,g)$. This class of metrics provides the setting for many interesting problems in geometric analysis and physics.

This class contains the `special' Poincar\'e--Einstein (PE) metrics, which by definition are the conformally compact metrics which satisfy
\begin{gather} \label{e1}
E := \Rc(g) + (n-1)g = 0,\qquad n = \dim M.
\end{gather}
There are differences in various parts of this theory depending on the parity of $n$, but in this paper we restrict ourselves
{\it entirely} to the case where $n$ is even. We emphasize our convention that $n = \dim M$ is the dimension of the bulk manifold,
which is in contrast to certain other papers on PE metrics where $n$ denotes the dimension of the conformal boundary instead.
It turns out that a choice of representative of the conformal infinity uniquely determines a boundary defining function and
a diffeomorphism of a collar neighbourhood of $\partial M$ with $\partial M \times (0,1)$. In terms of this identification,
there is a formal series solution $g$ to equation~\eqref{e1} called the Fefferman--Graham expansion. This expansion is
even up to order $n$ with respect to the special defining function, with coefficients locally determined by the choice of metric
in the conformal infinity. Furthermore, the first odd term in the expansion has vanishing trace. Using these properties,
one may then compute the volumes of compact truncations in the exhaustion of the manifold determined by these special
boundary defining functions. Discarding the divergent terms in the expansion of this one-parameter family of volumes
defines the renormalized volume. Considerations of parity in this expansion then shows that this renormalized volume
is independent of the choice of representative of the conformal infinity.

We study here this renormalized volume in a slightly more general setting, where the asymptotic expansion of the metric
has qualitatively similar features. The first possibility is to relax~\eqref{e1} and simply require that $g$ is asymptotically
Poincar\'e--Einstein (APE) in the sense that $|E|_g = O(\rho^n)$ for some (hence any) boundary defining function $\rho$. In
this case one has the same Fefferman--Graham expansion of the metric and the same volume renormalization scheme as
in the PE case. However, we relax these conditions even further, and require only that the expansion of the metric be even
to a critical order and that the first odd term have vanishing trace. We call AH metrics satisfying the first condition {\it
partially even}, and metrics satisfying both conditions {\it volume renormalizable}. Thus both PE metrics and APE metrics are
volume renormalizable. However, no special properties are required of the nonzero coefficients in the expansion of
a volume renormalizable metric, and hence the $E$ tensor need only vanish to second order.

AH metrics with various assumptions on the evenness of the expansion have been considered frequently before,
see \cite{Guillarmou, MM} and the more recent work by Vasy \cite{V} regarding the role of evenness in establishing
the meromorphic extension and properties of the resolvent of the Laplacian of an AH metric.

We also consider the normalized Ricci flow
\begin{gather} \label{normRF}
\partial_t g = -2 ( \Rc(g) + (n-1)g ), \qquad t \in [0,T), \qquad g(0) = \gin.
\end{gather}
The $(n-1) g$ term ensures that the conformal infinity is fixed in time. Observe that hyperbolic metrics and more generally
PE metrics are stationary points of this flow.

The preservation of the APE condition under the normalized Ricci flow and the variation of renormalized volume was studied
in our earlier paper~\cite{BMW}. Here we turn to these questions in the more general class of partially even and volume renormalizable
metrics. We now state our main results.

\begin{Theorem}\label{TheoremA} Suppose that $\big(M,\gin\big)$, $\dim M =: n = 2m$, is partially even, and let $g(t)$ be a~solution of~\eqref{normRF} with $g(0) = \gin$
with maximal interval of existence $[0,T_0)$. Then $(M, g(t))$ remains partially even for $t < T_0$. If, furthermore, $\big(M,\gin\big)$ is volume renormalizable then $(M, g(t))$ remains volume renormalizable for $t < T_0$.
\end{Theorem}
